\section{The terminal M\"obius layer and a sign-blind obstruction}
\label{sec:tpc34-terminal-barrier}

The orbit-sliced gate is a same-time row energy, and the unsigned
incidence estimates developed here do not by themselves control it.  On
the high--\(\beta\) endpoint packet this becomes explicit before any
affine reflection: the largest divisor of every target lies in the
ultra-divisor range, with a coefficient that may vanish when the target
is not squarefree.

\subsection{Exact terminal/proper-divisor splitting}

On the endpoint packet, \eqref{eq:tpc34-terminal-range-condition} gives
\(T<N_m(j)<U_0\).  Therefore
\begin{align}
 C_m(j)
 &=C_m^{\rm term}(j)+C_m^{\rm prop}(j),
 \label{eq:tpc34-terminal-split}\\
 C_m^{\rm term}(j)
 &:=
 a_R(N_m(j))
 =-\mu(N_m(j))\log N_m(j),
 \label{eq:tpc34-terminal-term}\\
 C_m^{\rm prop}(j)
 &:=
 \sum_{\substack{T<u<N_m(j)\\u\mid N_m(j)}}a_R(u).
 \label{eq:tpc34-proper-term}
\end{align}
This is an exact identity, including when \(\mu(N_m(j))=0\).

Insert \eqref{eq:tpc34-terminal-split} into
\eqref{eq:tpc34-left-orbit-slice} and write
\begin{equation}
 Y_{\gamma,j}^L
 =Y_{\gamma,j}^{L,\rm term}
 +Y_{\gamma,j}^{L,\rm prop},
 \label{eq:tpc34-Y-terminal-split}
\end{equation}
with corresponding squared energies \(V_L^{\rm term}\) and
\(V_L^{\rm prop}\); define the right versions similarly.  One always has
\begin{equation}
 V_L\le2(V_L^{\rm term}+V_L^{\rm prop}),
 \label{eq:tpc34-terminal-triangle}
\end{equation}
and similarly on the right.  This inequality gives a sufficient
layerwise route, but no converse lower bound: terminal and proper-divisor
vectors may cancel before squaring.

\subsection{The fused terminal sign}

Let \(m=\ell d\).  The row coefficient
\eqref{eq:tpc34-row-coefficients} contains the sign \(\mu(d)\), while
\eqref{eq:tpc34-terminal-term} contains
\(-\mu(\ell dj+h)\).  Target primitivity gives
\begin{equation}
 (d,\ell dj+h)=1.
 \label{eq:tpc34-row-target-coprime}
\end{equation}
Hence, whenever the two M\"obius factors are nonzero,
\begin{equation}
 \mu(d)\mu(\ell dj+h)
 =\mu\bigl(d(\ell dj+h)\bigr).
 \label{eq:tpc34-fused-terminal-sign}
\end{equation}

On the principal projector layer \(v=1\), the terminal divisor has
complementary factor \(s:=N_m(j)/u=1\).  Freezing \((\ell,j)\), the
resulting row sum has the schematic but literal sign pattern
\begin{equation}
 \sum_{d\asymp D}
 \mu(d)\mu(\ell jd+h)W_{\ell,j,\gamma}(d),
 \label{eq:tpc34-terminal-affine-Chowla}
\end{equation}
where \(W_{\ell,j,\gamma}\) retains the actual smooth, residue, and
opposite-row mask factors.  After restoring the \(\ell\)-sum and
expanding its same-time square, one obtains
\begin{equation}
 \mu(d_1)\mu(\ell_1d_1j+h)
 \mu(d_2)\mu(\ell_2d_2j+h).
 \label{eq:tpc34-four-Mobius-row}
\end{equation}
The averaging over \(j,\ell_1,\ell_2,d_1,d_2\) and the common opposite
row makes this weaker than a pointwise theorem for every affine
direction.  It is nevertheless sign-sensitive: the congruence and gcd
identities proved above do not remove
\eqref{eq:tpc34-four-Mobius-row}.

\begin{remark}[Relation to Chowla technology]
TPC-33 observes that the unrestricted fixed-direction family contains
\(\sum_n\mu(n)\mu(n+1)\), whereas logarithmically averaged fixed-form
results and shift-averaged results have different hypotheses
\citep{Tao2016LogChowla,MatomakiRadziwillTao2015,WangTPC33}.  We do not
claim that the physically averaged operator gate
\eqref{eq:tpc34-orbit-operator-gate} is equivalent to full two-point
Chowla.  We claim only that its terminal layer is a genuine averaged
affine M\"obius correlation and is not controlled by the unsigned
incidence estimates in this paper.
\end{remark}

\subsection{A sharp abstract majorant model}

The next finite example shows why closing the Gram-copy diagonal does not
close the full Gram energy.  It also records the precise limitation of a
sign-blind terminal majorant.

\begin{proposition}[Coherent zero-pair-diagonal model]
\label{prop:tpc34-coherent-majorant}
Let the input and opposite row sets both have \(M\ge2\) elements, let the
orbit have \(J_0\ge1\) times, and take
\begin{equation}
 \gamma_\alpha=1,\qquad
 \mathsf H_\gamma(j)=1,\qquad
 C_\alpha(j)=1,\qquad
 \mathfrak A_{\alpha,\gamma}(j)=\one_{\alpha\ne\gamma}.
 \label{eq:tpc34-coherent-model-data}
\end{equation}
Then
\begin{align}
 K_{\gamma,\alpha}
 &=J_0\one_{\alpha\ne\gamma},
 \label{eq:tpc34-coherent-model-K}\\
 E
 &=M(M-1)^2J_0^2,
 \label{eq:tpc34-coherent-model-E}\\
 E_\Delta
 &=M(M-1)J_0^2,
 \label{eq:tpc34-coherent-model-diagonal}\\
 E_{\ne}
 &=M(M-1)(M-2)J_0^2.
 \label{eq:tpc34-coherent-model-offdiagonal}
\end{align}
The orbit-sliced energy is
\begin{equation}
 V=M(M-1)^2J_0,
 \qquad E=J_0V.
 \label{eq:tpc34-coherent-model-V}
\end{equation}
\end{proposition}

\begin{proof}
For every opposite row \(\gamma\), exactly \(M-1\) input rows survive.
Thus \(B_\gamma=(M-1)J_0\), which proves
\eqref{eq:tpc34-coherent-model-E}.  There are \(M(M-1)\) nonzero matrix
entries, each of squared modulus \(J_0^2\), proving
\eqref{eq:tpc34-coherent-model-diagonal}; subtraction gives the
off-diagonal formula.  At every \((\gamma,j)\), the orbit-sliced row sum
is \(M-1\), proving \eqref{eq:tpc34-coherent-model-V}.
\end{proof}

For \(M\asymp Q\) and \(J_0\asymp J\), this model has
\[
 E\asymp Q^3J^2,\qquad V\asymp Q^3J.
\]
It respects the equal-pair deletion, the row cardinality, both row norm
scales, the orbit length, and all pointwise envelopes.  It may also be
viewed as retaining one abstract terminal atom per row-time pair and
replacing its fused sign by \(+1\).  It is \emph{not} asserted to be the
physical divisor kernel.

\begin{corollary}[Envelope-only method boundary]
\label{cor:tpc34-sign-blind-boundary}
No estimate that uses only the listed cardinalities, pointwise envelopes,
and equal-pair deletion can uniformly improve
\eqref{eq:tpc34-absolute-baseline} to \(O(Q^3)\) over the corresponding
abstract majorant class.
\end{corollary}

The corollary does not say that the physical energy is large and does not
exclude a stronger theorem using additional arithmetic incidence.  The
specific unsigned incidence bounds proved here save factors on selected
collision layers but remain above the complete gate.  A physical proof
may use the fused signs in
\eqref{eq:tpc34-fused-terminal-sign}, cancellation between terminal and
proper-divisor layers in \eqref{eq:tpc34-Y-terminal-split}, or another
sign-sensitive spectral mechanism.  The listed sign-blind data alone do
not supply the missing powers.  This is a local
method obstruction consistent with the classical parity boundary
\citep{HeathBrown1982Parity}; it is not a parity-breaking theorem.
